\scanpage{160}\mspage{1\unsure{5|6}}
$G\longmapsto\Delta G=\mathbb{R}\underline{\mathrm{Hom}}\,(G,\mathbb{Q}/\mathbb{Z})$ : on \unsure{l'appelle|l'appellera} encore dualité de Serre (``cohomologique''). Ce sera une magnifique autodualité, si on se place dans une catégorie dérivée où les $\underline{H}^{i}$ des complexes envisagés sont (à des \unsure{passages|p.n.g} \unsure{à|;} la limite près) des extensions de groupes étales par des groupes unipotents connexes. Or on ne rencontre que de tels complexes, en ``intégrant'' des coeff. finis $F$ sur $X$ par $Rf_!$ ou $Rf_*$. Ceci posé, passant à la limite sur $n$ dans la formulation initiale (ou encore partant avec $(\mathbb{Q}/\mathbb{Z})_k$ sur $k$, au lieu des $(\mathbb{Z}/n\mathbb{Z})_k$, et formant $f^!(\mathbb{Q}/\mathbb{Z})_k=(\mu_\infty)_X$), la formule de dualité prend la forme
\diagram{\[
\begin{tikzcd}[ampersand replacement=\&,column sep=0.8em,row sep=0.1em]
\& \text{``shift'' des dimensions} \\
\boxed{\Delta(Rf_!(F))\simeq Rf_*(DF[2])} \arrow[ur,start anchor=east,end anchor=west,shorten <=-0.8em] \&
\end{tikzcd}
\]}
\unsure{où} $D$ est la ``dualité de Cartier'' $R\underline{\mathrm{Hom}}\,(F,\mu_\infty)$ (ou $R\underline{\mathrm{Hom}}\,(F,\mathbb{G}_m)$ au choix ?), et $\Delta$ la dualité de Serre : cohomologie à support propre et à support quelconque s'échangent\inserted{ par dualité}, quand on prend ``en haut la dualité de \unsure{Cartier|Cartin}, ``en bas'' la dualité de Serre.

La validité de la formule de dualité ne fait aucun doute -- le travail principal pour
% archive sheet number: 160
